\documentclass[10pt]{beamer}


\usepackage[utf8]{inputenc}%pacchetto lettere accentate in input
\usepackage[english]{babel}%pacchetto lingue
\usepackage{amsmath,mathtools,amssymb,amsfonts,braket,bbm}%pacchetti font matematici, simboli matematici e parentesi
\usepackage{graphicx,caption,subfig}
\usepackage{enumitem} %utile alla gestione e personalizzazione degli elenchi
\usepackage{mathrsfs,amsthm} %utile per la gestione dei comandi newtheorem, \begin{proof}...
\usepackage[autostyle,italian=guillemets]{csquotes} %parole carine titolo
\usepackage{tikz} %pacchetto per disegnare, sufficienti librerie interne, nozioni su LatexPedia
\usetikzlibrary{arrows,automata,topaths,trees} %librerie di tikz per frecce, cammini, alberi e disegni di macchine
\usetikzlibrary{decorations.markings} %per decorare frecce, in particolare per fare la punta al centro della freccia
\usetikzlibrary{decorations.pathreplacing,calligraphy} %per le parentesi graffe nei disegni
%\usepackage{xr-hyper} 
\usetikzlibrary { positioning }
\usepackage{hyperref}% riferimenti cliccabili
%\externaldocument[D-]{tesi}[tesi.pdf]
\hypersetup{linktoc=all,colorlinks,
	citecolor= black,
    filecolor=black,
    linkcolor=black,
    urlcolor=black}
    %definizione dei caratteri raffinati: usati per le lettere di probabilità, valor medio e per gli insiemi dei naturali, reali, interi
\def\CC {\mathbb{C}}
\def\EE {\mathbb{E}}
\def\NN {\mathbb{N}}
\def\PP {\mathbb{P}}
\def\RR {\mathbb{R}}
\def\ZZ {\mathbb{Z}}
\def\Earr {\Vec{E}}
%definizione dei caratteri calligrafici: usati per sigma-algebre, resistenze e conduttanze effettive
\def\R  {\mathcal{R}}
\def\C  {\mathcal{C}}
\def\E  {\mathcal{E}}
\def\F  {\mathcal{F}}
% definizione funzione indicatrice
\def\ind {\mathbbm{1}}
% definizione comando: resistenza effettiva
\newcommand{\Reff}[2]{\R_{\text{eff}}(#1 \leftrightarrow #2) }
% definizione comando: freccetta vettore ma da destra a sinistra
\newcommand{\cev}[1]{\reflectbox{\ensuremath{\vec{\reflectbox{\ensuremath{#1}}}}}}


%definizioni di Definizione,Teorema, Claim, Corollario (del teorema), Corollario (del Claim), Esempio
\theoremstyle{definition}
\newtheorem{definizione}{Definizione}

\theoremstyle{definition}
\newtheorem{esempio}[definizione]{Esempio}

\theoremstyle{definition}
\newtheorem{controesempio}[definizione]{Controesempio}

\theoremstyle{plain}
\newtheorem{proposizione}[definizione]{Proposizione}

\theoremstyle{plain}
\newtheorem{corollario}[definizione]{Corollario}

\theoremstyle{plain}
\newtheorem{teorema}[definizione]{Teorema}

\theoremstyle{definition}
\newtheorem{osservazione}[definizione]{Osservazione}


\usetheme{Padova}

\title{\scshape{Inverted Pendulum}}
\subtitle{\fontsize{12}{12} \scshape{Model Predictive Control}\vspace{.5cm}}
\author{\fontsize{9}{9} \scshape{Student: Lorenzo Trubian} \\ \hspace{40pt} \fontsize{7}{7} \scshape{n° 2125512}\vspace{.5cm}}
\date{\fontsize{8}{8} \scshape{2 January 2026}}

\parskip 6pt
\begin{document}

% TITOLO
	\maketitle
% SOMMARIO
%	\begin{frame}{Sommario}
%		\tableofcontents
%	\end{frame}
% INTRODUZIONE 
%	\section{Introduzione}
%
%	\begin{frame}{Introduzione}
%		Ci proponiamo di descrivere alcune proprietà delle passeggiate aleatorie su grafi avvalendoci del linguaggio dei circuiti elettrici. L'interesse principale è quello di determinare se un grafo connesso dotato di pesi positivi è ricorrente o transiente e di applicare i risultati ottenuti ai casi particolari di $\ZZ^2$ e $\ZZ^3$. \vspace{.5em}
%		\begin{itemize}
%			\item   \vspace{.5em}
%			\item   \vspace{.5em}
%			\item 
%		\end{itemize}
%	\end{frame}

% SEZIONE FUNZIONI ARMONICHE E PASSEGGIATE ALEATORIE
	\section{Extended model}
	\begin{frame}{Catene di Markov e Reti}

			\begin{block}{Definizione (Rete)}
				Una \textbf{rete}
			\end{block}
			\vspace{-.5cm}
		\begin{figure}
			\caption{Esempio di rete. In ogni vertice è segnato $\pi(x)$.}


		\end{figure}
	\end{frame}
	
	\begin{frame}{Catene di Markov e Reti}
			Identificheremo in maniera univoca una passeggiata aleatoria su una rete ponendo le probabilità di transizione della catena di Markov $\{X_n\}_{n\geq0}$
			\begin{equation}
			p_{xy} := \PP(X_{t+1} = y \mid X_t = x) = \frac{c_{xy}}{\pi(x)}.
			\end{equation}
			Viceversa, ad ogni catena di Markov reversibile e omogenea si può associare una rete $G=(V,E)$ ponendo in corrispondenza i vertici con gli stati della catena, congiungendo ogni coppia di vertici $x,y$ qualora $p_{xy}>0$ e infine assegnando all'arco $(x,y)$ il peso
			\begin{equation}
			c_{xy}:=\pi(x)p_{xy}.
			\end{equation}
			Poiché la catena è reversibile, il peso è ben definito. Infatti $c_{xy} = \pi(x)p_{xy} = \pi(y)p_{yx} = c_{yx} $.
	\end{frame}
	\begin{frame}{Catene di Markov e Reti}
	Precisiamo la seguente notazione:
	\begin{enumerate}[label=-]
	\item Useremo $\PP_x(\cdot)$ per indicare la probabilità condizionata $\PP(\cdot\mid X_0 =x)$.\\
	\item Indicheremo il \textbf{tempo di visita} di un sottoinsieme di stati $A$ con
	\[
	\tau_A = \inf\Set{n\geq0\mid X_n \in A}.
	\]
	\item Infine, indicheremo con 
	\[
	\tau_A^+ = \inf\Set{n\geq1\mid X_n\in A}.
	\]
	\end{enumerate}
	Se $X_0 \notin A$ allora $\tau_A^+ = \tau_A$ quasi certamente.\par
	In modo colloquiale, se $X_0 \in A$ il tempo $\tau_A^+$ può essere inteso come ``tempo di ritorno''.
	\end{frame}
	
	
\section{Funzioni armoniche e Passeggiate aleatorie}
	\begin{frame}{Funzioni armoniche e Passeggiate aleatorie}
		\begin{block}{Definizione}
		Una funzione $h:V \rightarrow \RR $ è \textbf{armonica} su un vertice $x\in V$ se
		\begin{equation}
			h(x) = \frac{1}{\pi(x)}\sum_{y:y \sim x} c_{xy}h(y).
		\end{equation}
		Se $h$ è armonica su ogni elemento di un insieme $W$, allora diremo che $h$ è armonica 				su $W$.
	\end{block}
		\begin{block}{Definizione}
		Fissata una rete $G=(V,E)$ e due suoi vertici distinti $a,z \in V$, diremo \textbf{voltaggio} una funzione $g$ armonica su $V\setminus\set{a,z}$. Diremo che un voltaggio è unitario se $g(z) = 1$ e $g(a) = 0$.
		\end{block}
	\end{frame}


	\begin{frame}{Funzioni armoniche e Passeggiate aleatorie}
		\begin{alertblock}{Proposizione}
				\label{claim:esistenza-rete-finita}
				Data una rete finita $G=(V,E)$, fissati due vertici $a\neq z \in V$, la funzione \[
				g(x) := \PP_x(\tau_z<\tau_a).
				\]
				è l'unico voltaggio unitario.
		\end{alertblock}
		\vspace{-.5cm}
		\begin{figure}
			\caption{Rete con pesi positivi sugli archi e valori del voltaggio unitario sui vertici $a$ e $z$.}
		\end{figure}
	\end{frame}
	\section{Flussi e Correnti}
	\begin{frame}{Flussi e Correnti}
	Per un grafo $G=(V,E)$, denoteremo con $\Earr$ l'insieme di tutti archi orientati di $G$. In altre parole, $(x,y) \in \Earr$ se e solo se $\Set{x,y} \in E$ e in tal caso si ha anche che $(y,x)\in \Earr$.
		\begin{block}{Definizione (Flusso)}
			Un \textbf{flusso da $a$ a $z$} in una rete $G$ è una funzione $\theta:\Earr\rightarrow\RR$ che soddisfa:
			\begin{enumerate}[label=(\roman*)]
			\item per ogni $\Set{x,y}\in E$ si ha $\theta(xy) = - \theta(yx)$ (\textbf{antisimmetria}),
			\item per ogni $x\notin \Set{a,z}$ si ha $\sum_{y:y\sim x} \theta(xy) = 0$ (\textbf{legge dei nodi di Kirchhoff}).
			\end{enumerate}
		\end{block}
	\end{frame}
	\begin{frame}{Flussi e Correnti}

		\begin{block}{Definizione (Flusso associato ad un voltaggio)}
			Dato un voltaggio $h$, il \textbf{flusso di corrente} $\theta = \theta_h$ associato ad $h$ è definito da 
			\begin{equation}
			\label{eq:legge-ohm}
			\theta(xy) = c_{xy}(h(y)-h(x)).
			\end{equation}
		\end{block}
		\vspace{-0.3cm}
		\footnotesize La relazione (\ref{eq:legge-ohm}) è nota come \emph{legge di Ohm}: $V = R I$.
		\normalsize
			\begin{block}{Definizione (Corrente)}
			La \textbf{corrente} di un flusso $\theta$ da $a$ a $z$ è definita come
			\begin{equation}
			\|\theta\| := \sum_{x:x\sim a} \theta(ax).
			\end{equation}
		\end{block}
		\footnotesize{Si noti che la corrente di un flusso, benché rappresentata con una notazione di solito usata per le norme, non è in alcun modo una norma. Essa infatti può assumere valori negativi.}
	\end{frame}
	
\section{Resistenza effettiva}
	\begin{frame}{Resistenza effettiva}
		\begin{alertblock}{Proposizione}
			Fissata una rete $G=(V,E)$ e due suoi vertici $a,z$, per ogni voltaggio non costante $h$ e flusso di corrente $\theta$ ad esso associato, il rapporto
			\[
			\frac{h(z) - h(a)}{\|\theta\|}
			\]
			è una costante positiva che non dipende da $h$.
		\end{alertblock}
				\begin{block}{Definizione}
			Fissata una rete $G$ e dato un voltaggio non costante e il flusso di corrente associato, diremo \textbf{resistenza effettiva} tra $a$ e $z$ il rapporto
			\begin{equation}
			\R_{\text{eff}}(a\leftrightarrow z; G) := \frac{h(z)-h(a)}{\|\theta\|}.
			\end{equation}
		\end{block}
	\end{frame}
	\begin{frame}{Resistenza effettiva}
		\begin{alertblock}{Teorema}
			Per ogni rete finita si ha:
			\begin{equation}
			\label{eq:resistenza-effettiva-passeggiate}
			\R_{\text{eff}}(a\leftrightarrow z; G)=\frac{1}{\pi(a)\PP_a(\tau_z<\tau_a^+)}.
			\end{equation}
		\end{alertblock}
		La dimostrazione si serve del voltaggio unitario $h(x) = \PP_x(\tau_z<\tau_a)$, della legge delle probabilità totali e in particolar modo della proprietà di Markov.
	\end{frame}
	\section{Energia e Principio di Thomson}
	\begin{frame}{Energia e Principio di Thomson}
	\begin{block}{Definizione}
		L'\textbf{energia} di un flusso $\theta$ da $a$ a $z$, denotata con $\E(\theta)$ è così definita:
		\[
		\E(\theta):=\frac{1}{2}\sum_{\vec{e}\in\vec{E}}r_{\vec{e}}\theta(\vec{e})^2 = \sum_{e\in E} r_e\theta(e)^2.
		\]
	\end{block}
	\vspace{-.45cm}
	\footnotesize Si osservi che la seconda sommatoria avviene sull'insieme degli archi non orientati, ma dal momento che $\theta(xy)^2 = \theta(yx)^2$, la definizione è ben posta. Inoltre, l'energia è sempre maggiore o uguale a zero.
	\normalsize
	\begin{alertblock}{Teorema (Principio di Thomson)}
	\label{principio}
		Data una rete finita $G$ si ha
		\[
		\Reff{a}{z;G} = \inf\Set{\E(\theta) \mid \|\theta\| =1, \;\theta\text{ è un flusso da $a$ in $z$}},
		\]
		e l'unico elemento per il quale il minimo è raggiunto è il flusso di corrente unitaria.
	\end{alertblock}
	\end{frame}
	\section{Grafi infiniti}
	\begin{frame}{Grafi infiniti}
		Sia $G=(V,E)$ un grafo connesso e infinito dotato di conduttanze $\{c_e\}_{e\in E}$. Assumeremo sempre che la rete sia \textbf{localmente finita} ovvero che per ogni vertice $x\in V$ si ha $\sum_{y:y\sim x} c_{xy} <\infty$. 
		\begin{block}{Definizione}
		Sia $\{G_n\}$ una successione di sottografi finiti e connessi di $G$ tale per cui $\bigcup_{n\in\NN}G_n=G$ e $G_n\subset G_{n+1}$; diremo che una tale successione è \textbf{esaustiva} di $G$. Inoltre, identificheremo tutti i vertici di $G\setminus G_n$ con un singolo vertice $z_n$.
		\end{block}

		\begin{alertblock}{Proposizione}
			Data una successione esaustiva $\{G_n\}$ di $G$, il limite
			\begin{equation}
			\lim_{n\rightarrow\infty} \Reff{a}{z_n;G_n\cup\Set{z_n}} 
			\end{equation}
			esiste sempre e non dipende dalla scelta della successione esaustiva.
		\end{alertblock}
	\end{frame}
	\begin{frame}{Grafi infiniti}
		\begin{block}{Definizione}
			In una rete infinita, la resistenza effettiva da un vertice $a$ a $\infty$ è
			\[
			\Reff{a}{\infty}:=	\lim_{n\rightarrow\infty} \Reff{a}{z_n;G_n\cup\Set{z_n}}.
			\]
		\end{block}
		\begin{block}{Definizione}
			Una rete $(G,\{r_e\}_{e\in E})$ è detta \textbf{ricorrente} se $\PP_a(\tau_a^+=\infty)=0$, ovvero se la probabilità che una passeggiata aleatoria che inizi in $a$ non torni mai in $a$ è zero. Altrimenti, la rete è detta \textbf{transiente}.
		\end{block}
		\vspace{-.3cm}
		La relazione tra resistenza effettiva e catene di Markov (\ref{eq:resistenza-effettiva-passeggiate}) si estende ai grafi infiniti come segue:
		\begin{equation}
		\Reff{a}{\infty} = \frac{1}{\pi(a)\PP_a(\tau_a^+ = \infty)},
		\end{equation}
		con la convenzione che $1/0=\infty$.
	\end{frame}
	
	\begin{frame}{Grafi infiniti}
		\begin{alertblock}{Teorema (Principio di Thomson per reti infinite)}
			Sia $G$ una rete infinita, allora
			\[
			\Reff{a}{\infty} = \inf\Set{\E(\theta)\mid\theta\text{ è un flusso da $a$ a $\infty$ di corrente unitaria}}.
			\]
		\end{alertblock}

			\begin{alertblock}{Corollario}
				Sia $G$ una rete infinita. Le seguenti affermazioni sono equivalenti:
				\begin{enumerate}[label=(\roman*)]
				\item $G$ è transiente.
				\item Esiste un vertice $a\in V$ tale che $\Reff{a}{\infty}<\infty$. Di conseguenza, tutti i vertici soddisfano questa proprietà.
				\item Esiste un vertice $a\in V$ e un flusso unitario $\theta$  da $a$ a $\infty$ con $\E(\theta) <\infty$. Di conseguenza, tutti i vertici soddisfano questa proprietà.
			\end{enumerate}
		\end{alertblock}
	\end{frame}
	\begin{frame}{Grafi infiniti}
		\begin{block}{Definizione (Taglio)}
			Un \textbf{taglio} $\Gamma \subseteq E(G)$ che separa $a$ da $z$ è un insieme di archi tale per cui ogni cammino da $a$ a $z$ possiede un arco in $\Gamma$. In una rete infinita $G=(V,E)$, diremo che $\Gamma\subset E$ è un taglio che separa $a$ da $\infty$ se ogni cammino semplice e infinito che inizi in $a$ possiede un arco in $\Gamma$.
		\end{block}	
	\end{frame}
	\begin{frame}{Grafi infiniti}
		\begin{alertblock}{Teorema (Disuguaglianza di Nash-Williams)}
			Data una rete finita, siano $\{\Gamma_n\}$ tagli a due a due disgiunti che separano $a$ da $z$. Allora si ha
			\begin{equation}
			\Reff{a}{z}\geq\sum_{n}\left(\sum_{e\in\Gamma_n}c_e\right)^{-1}.
			\end{equation}
		\end{alertblock}
		\begin{alertblock}{Corollario}
			In una rete infinita, se esiste una collezione $\{\Gamma_n\}$ di tagli a due a due disgiunti che separato $a$ da $\infty$ tali per cui
			\[
			\sum_n\left(\sum_{e\in\Gamma_n}c_e\right)^{-1} = \infty
			\]
	allora la rete è ricorrente.
		\end{alertblock}
	\end{frame}
	
	\begin{frame}{Grafi infiniti: $\ZZ^2$}
		\begin{alertblock}{Teorema}
		$\ZZ^2$ dotato di pesi unitari è ricorrente
		\end{alertblock}
		\vspace{-.3cm}
			\emph{Dimostrazione.} Definiamo $\Gamma_n$ come l'insieme degli archi $\{(x,y),(x,y+1)\}$ con $|x|\leq n$ e $\min\Set{|y|,|y+1|}=n$ e di quelli della forma $\{(x,y),(x+1,y)\}$ con $|y|\leq n$ e $\min\Set{|x|,|x+1|}=n$. Allora, $\{\Gamma_n\}_{n\geq1}$ è una collezione di tagli a due a due disgiunti che separano $0$ da $\infty$. Inoltre, $|\Gamma_n|=4(2n+1)$ e quindi \footnotesize$ \sum_n\left(\sum_{e\in\Gamma_n}c_e\right)^{-1} = \sum_n\frac{1}{|\Gamma_n|} = \infty$. \qed
		\vspace{-.4cm}
		\begin{figure}
			\caption{Alcuni sottografi di $\ZZ^2$. Sono rappresentati i tagli $\Gamma_1$ e $\Gamma_3$ (in azzurro) e $\Gamma_2$ (in rosso). L'origine è cerchiata.}
			\vspace{-.8cm}
			\begin{minipage}[t][3cm][c]{0.3\textwidth}
			\end{minipage}
			\begin{minipage}[t][3cm][c]{0.3\textwidth}
			\end{minipage}
			\begin{minipage}[t][3cm][c]{0.3\textwidth}
			\end{minipage}
		\end{figure}
	\end{frame}
	
	\begin{frame}{Grafi infiniti}
	\begin{alertblock}{Proposizione}
		Considerato un cammino $\gamma$ da $a$ in $z$, sia definito
		\begin{align*}
		\theta_\gamma(\vec{e}) = &~ (\text{numero di volte } \vec{e} \text{ è attraversato da } \gamma ) \\ 
		& -(\text{numero di volte } \cev{e} \text{ è attraversato da } \gamma),
		\end{align*}
		dove $\vec{e},\cev{e}$ indicano i due orientamenti di un arco $e$ di $G$. Poniamo
		\[
		\theta(\vec{e}) = \EE(\theta_\gamma(\vec{e})).
		\]
		Allora $\theta$ è un flusso da $a$ in $z$ con $\|\theta\|=1$.
	\end{alertblock}
	\end{frame}
	
	\begin{frame}[allowframebreaks]{Grafi infiniti: $\ZZ^3$}
	\vspace{.05cm}
	\begin{alertblock}{Teorema}
	$\ZZ^3$ dotato di pesi unitari è transiente.
	\end{alertblock}
	\emph{Dimostrazione.} Per $R>0$ denotiamo con $B_R = \Set{(x,y,z): x^2+y^2+z^2 \leq R^2}$ la palla di raggio $R$ in $\RR^3$. Posto $V_R = B_R \cap \ZZ^3$ e sia $\partial V_R$ l'insieme dei vertici esterni a $V_R$ e confinanti con $V_R$, ovvero l'insieme dei vertici non appartenenti a $V_R$ che sono estremi di archi il cui altro estremo è in $V_R$. \par
	Costruiamo ora una passeggiata aleatoria $\gamma$ da $0$ a $\partial V_R$ scegliendo casualmente un punto $p$ in $\partial B_R = \Set{(x,y,z):x^2+y^2+z^2=R}$, tracciando una linea retta da $0$ a $p$ in $\RR^3$, considerando l'insieme di punti in $\RR^3$ distanti dalla linea retta al più $10$ e scegliendo in modo del tutto arbitrario un cammino in $\ZZ^3$ contenuto in quell'insieme. La costante $10$ è scelta arbitrariamente in modo da garantire l'esistenza di un cammino discreto per ogni punto $p\in\partial B_R$. \par
	Per la Proposizione precedente la funzione $\theta(\vec{e}) = \EE(\theta_\gamma(\vec{e}))$ corrisponde ad un flusso da $0$ a $\partial V_R$, dove $\gamma$ sono i cammini costruiti come appena descritto. Per stimare l'energia di questo flusso, notiamo che se $\vec{e}$ è un arco distante $r\leq R$ dall'origine, allora la probabilità che sia attraversato da un cammino $\gamma$ è $O(r^{-2})$. Inoltre, ci sono $O(r^2)$ di questo tipo di archi. Quindi, l'energia di un flusso è al più
	\begin{equation}
	\E = O\left(\sum_{r=1}^Rr^2\cdot(r^{-2})^2\right)\leq C,
	\end{equation}
	per una qualche costante $C<\infty$ che non dipende da $R$. Per il Principio di Thomson deduciamo che $\Reff{0}{\partial V_R}\leq C$ per ogni $R$, e di conseguenza che $\ZZ^3$ è transiente. \qed
	\end{frame}
	\begin{frame}
	\fontsize{12}{12} \scshape{Grazie dell'attenzione!}\vspace{.5cm}
	\end{frame}
\end{document}
